\clearpage\thispagestyle{empty}
\pdfbookmark[0]{Quelle zu Ceres, Elemente XII}{ceres1808xii-sources}
\begin{center}\textbf{Quelle zu Ceres, Elemente XII}\end{center}
\noindent Diese KI-Übertragung ins heutige Deutsch folgt der Anzeige vom 10. März 1808 in \href{https://gdz.sub.uni-goettingen.de/id/PPN236007467}{Gauß’ \textit{Werken}, Bd.~VI} (Göttingen, 1874), S.~299--300. Die Anzeige erschien unter der Überschrift \textit{Göttingen} in den \textit{Göttingischen gelehrten Anzeigen}, Stück~40, \href{https://manifests.sub.uni-goettingen.de/iiif/image/gdz:PPN319721507_1808:00000397/full/1600,/0/default.jpg}{S.~393}, \href{https://manifests.sub.uni-goettingen.de/iiif/image/gdz:PPN319721507_1808:00000398/full/1600,/0/default.jpg}{394} und \href{https://manifests.sub.uni-goettingen.de/iiif/image/gdz:PPN319721507_1808:00000399/full/1600,/0/default.jpg}{395}. Der Text spricht von Gauß in der dritten Person.
\par\medskip
Der Erstdruck vergleicht fünf Oppositionen. Die Werkausgabe enthält außerdem den Vergleich für 1808 und die Unterschiede für die Elemente XIII. Schering weist darauf hin, dass Gauß beides handschriftlich hinzugefügt hat. Auch der Hinweis auf Harding in eckigen Klammern steht in der Werkausgabe. Das leere Feld für die Breite von 1808 und das unregelmäßige Zeichen in der Punktreihe zur wahren Länge folgen der Werkausgabe. KI-Übertragung und KI-Satz.
\clearpage
